\section*{Hoofdstuk \ref{ch:b1:continuity} --- Limieten en continuïteit}

\begin{solution}{exo:b1:continuity:1}
Neem $u_n = \frac{1}{2\pi n + \pi/2}$ en $v_n = \frac{1}{2\pi n}$:
beide naderen tot $0^+$, terwijl $\sin\frac{1}{u_n} = 1$ en
$\sin\frac{1}{v_n} = 0$. Twee rijen, twee verschillende limieten van
de beelden: volgens \cref{thm:b1:continuity:seqchar} bestaat er geen
limiet in $0^+$.

$x \sin\frac1x$: ingeklemd door $\abs{x\sin\frac1x} \leq \abs x \to
0$, dus de limiet in $0$ bestaat en is gelijk aan $0$.
\end{solution}

\begin{solution}{exo:b1:continuity:2}
$f = \lfloor\cdot\rfloor$ is \omterm{def:b1:continuity:continuous}{continu} op $\R \setminus \Z$ (lokaal
constant) en discontinu in elke $n \in \Z$: linkerlimiet $n - 1$,
waarde $n$.

$g(x) = x - \lfloor x\rfloor$: dezelfde discontinuïteitspunten (de
identiteit is \omterm{def:b1:continuity:continuous}{continu}, dus $g$ erft de sprongen van $f$); in $n \in
\Z$ is de linkerlimiet $1 \neq 0 = g(n)$.

$h(x) = \lfloor x\rfloor + (x - \lfloor x\rfloor)^2$: op $\intco{n}{n
+ 1}$ geldt $h(x) = n + (x - n)^2$, daar \omterm{def:b1:continuity:continuous}{continu}; in $x = n$ is de
linkerlimiet $(n-1) + 1 = n = h(n)$: de sprongen heffen elkaar op.
$h$ is \omterm{def:b1:continuity:continuous}{continu} op $\R$ (en strikt stijgend).
\end{solution}

\begin{solution}{exo:b1:continuity:3}
$P(x) = x^5 - 3x + 1$: $P(-2) = -32 + 6 + 1 = -25 < 0$; $P(0) = 1 >
0$; $P(1) = -1 < 0$; $P(2) = 27 > 0$. Drie tekenwisselingen op de
disjuncte segmenten $\intcc{-2}{0}$, $\intcc{0}{1}$, $\intcc{1}{2}$:
volgens \cref{thm:b1:continuity:ivt} zijn er minstens drie
nulpunten. (Omdat de graad $5$ is, zijn er hoogstens vijf; een
studie van de variaties zou aantonen dat het er precies drie zijn.)
\end{solution}

\begin{solution}{exo:b1:continuity:4}
Zij $P = a_{2m+1}X^{2m+1} + \dots$ met $a_{2m+1} > 0$ (vervang
anders $P$ door $-P$). Door de dominante term af te zonderen krijgen
we $P(x) = a_{2m+1} x^{2m+1}\bigl(1 + o(1)\bigr)$ als $x \to
\pm\infty$: dus $P(x) \to +\infty$ in $+\infty$ en $-\infty$ in
$-\infty$. Kies $a$ met $P(a) < 0$ en $b$ met $P(b) > 0$: de
tussenwaardestelling op $\intcc{a}{b}$ levert een nulpunt.
\end{solution}

\begin{solution}{exo:b1:continuity:5}
Zij $g(x) = f(x) - x$, \omterm{def:b1:continuity:continuous}{continu} op $\intcc{0}{1}$. Omdat $f$ afbeeldt
in $\intcc{0}{1}$: $g(0) = f(0) \geq 0$ en $g(1) = f(1) - 1 \leq 0$.
Volgens \cref{thm:b1:continuity:ivt} is $g(c) = 0$ voor zekere $c$:
een vast punt.

De hypothesen zijn onmisbaar: op $\intcc{0}{1}$ heeft de discontinue
\omterm{def:b1:logic:map}{afbeelding} $f(x) = 1$ voor $x \leq \frac12$, $f(x) = 0$ voor $x >
\frac12$ geen vast punt; op het \omterm{prop:b1:reals:intervals}{interval} $\intoo{0}{1}$ (geen
segment) is $f(x) = \frac x2$ \omterm{def:b1:continuity:continuous}{continu} met waarden in $\intoo{0}{1}$
en zonder vast punt (de kandidaat $0$ ontbreekt); op $\R$ voldoet
$f(x) = x + 1$.
\end{solution}

\begin{solution}{exo:b1:continuity:6}
$g(x) = \cos x - x$ is \omterm{def:b1:continuity:continuous}{continu}, $g(0) = 1 > 0$, $g(1) = \cos 1 - 1 <
0$: er bestaat een oplossing in $\intoo{0}{1}$
(\cref{thm:b1:continuity:ivt}). Eenduidigheid: $g$ is strikt dalend
op $\R$ --- voor $x \leq 0$ is $g(x) \geq 1 - x > 0$, dus daar is er
sowieso geen nulpunt; en $g'(x) = -\sin x - 1 \leq 0$ met gelijkheid
alleen in geïsoleerde punten ($x \equiv -\frac\pi2 \bmod 2\pi$), dus
$g$ is strikt dalend (\cref{ch:b1:derivative}; anders gezegd: op
$\intcc{0}{1}$ is $\cos$ strikt dalend en $-x$ ook, dus $g$
eveneens). Een strikt monotone functie heeft hoogstens één nulpunt.
\end{solution}

\begin{solution}{exo:b1:continuity:7}
Leg $M = f(0) + 1$ vast. Er is een $A > 0$ met $f(x) \geq M$ voor
$\abs x \geq A$ (definitie van de twee oneindige limieten; neem de
grootste van beide drempels). Op het segment $\intcc{-A}{A}$ levert
de extremumstelling (\cref{thm:b1:continuity:evt}) een $c$ met $f(c)
= \inf_{\intcc{-A}{A}} f \leq f(0)$. Voor $\abs x \geq A$ geldt
$f(x) \geq M > f(0) \geq f(c)$. Dus is $f(c)$ het globale minimum.
\end{solution}

\begin{solution}{exo:b1:continuity:8}
Op het segment $\intcc{0}{T}$ is $f$ begrensd en bereikt het zijn
grenzen (\cref{thm:b1:continuity:evt}); wegens de periodiciteit zijn
dat de grenzen op heel $\R$, en ze worden nog steeds bereikt.

Zij $g(x) = f(x + \frac T2) - f(x)$, \omterm{def:b1:continuity:continuous}{continu}. Dan is
\[
g(0) + g\bigl(\tfrac T2\bigr)
= \bigl(f(\tfrac T2) - f(0)\bigr) + \bigl(f(T) - f(\tfrac T2)\bigr)
= f(T) - f(0) = 0 :
\]
$g(0)$ en $g(\frac T2)$ hebben tegengesteld teken (of een van beide
is nul), dus geeft de tussenwaardestelling op $\intcc{0}{T/2}$ een
$c$ met $g(c) = 0$, dat wil zeggen $f(c + \frac T2) = f(c)$.
\end{solution}

\begin{solution}{exo:b1:continuity:9}
Eerst de ongelijkheid: voor $0 \leq y \leq x$ is
\[
\bigl(\sqrt y + \sqrt{x - y}\bigr)^2 = x + 2\sqrt{y(x-y)} \geq x,
\]
dus $\sqrt x \leq \sqrt y + \sqrt{x - y}$, dat wil zeggen $\sqrt x -
\sqrt y \leq \sqrt{x - y}$. Bijgevolg geldt $\abs{\sqrt x - \sqrt y}
\leq \sqrt{\abs{x - y}}$ voor alle $x, y \geq 0$.

\omterm{def:b1:continuity:uniform}{Uniforme continuïteit}: zij $\varepsilon > 0$ gegeven; neem $\delta =
\varepsilon^2$; dan impliceert $\abs{x - y} \leq \delta$ dat
$\abs{\sqrt x - \sqrt y} \leq \sqrt\delta = \varepsilon$.

Niet lipschitziaans bij $0$: $\frac{\sqrt x - \sqrt 0}{x - 0} =
\frac{1}{\sqrt x} \to +\infty$ als $x \to 0^+$, dus geen enkele
constante $k$ kan alle differentiequotiënten domineren.
\end{solution}

\begin{solution}{exo:b1:continuity:10}
Stel dat $f$ \omterm{def:b1:logic:inj}{injectief} en \omterm{def:b1:continuity:continuous}{continu} is, maar niet strikt monotoon. Dan
zijn er $a < b < c$ in $I$ waarvoor $f(b)$ niet tussen $f(a)$ en
$f(c)$ ligt --- immers, als voor \emph{elk} drietal de middelste
waarde tussen de buitenste lag, zou $f$ monotoon zijn (vergelijk
twee willekeurige paren; een korte gevalsonderscheiding). Zeg $f(b)
> \max(f(a), f(c))$ (het andere geval is symmetrisch: vervang $f$
door $-f$). Kies $v$ met $\max(f(a), f(c)) < v < f(b)$. Door de
tussenwaardestelling toe te passen op $\intcc{a}{b}$ en op
$\intcc{b}{c}$ vinden we $u_1 \in \intoo{a}{b}$ en $u_2 \in
\intoo{b}{c}$ met $f(u_1) = v = f(u_2)$: twee verschillende punten
met hetzelfde beeld, in tegenspraak met de injectiviteit.
\end{solution}

\begin{solution}{exo:b1:continuity:11}
Uit $f(0) = f(0+0) = 2f(0)$ volgt $f(0) = 0$; uit $0 = f(x - x)$
volgt $f(-x) = -f(x)$. Stel $\alpha = f(1)$. Met inductie: $f(n) =
n\alpha$ voor $n \in \N$, en dan voor $n \in \Z$ wegens de
oneven\-heid. Voor $q \in \N^*$ geldt $q\,f(\frac pq) = f(p) =
p\alpha$ (tel $\frac pq$ $q$ keer bij zichzelf op), dus $f(\frac pq)
= \alpha\frac pq$: op $\Q$ is $f = \alpha\,\mathrm{id}$.

Zij nu $x \in \R$ en $(r_n)$ een rij rationale getallen met $r_n \to
x$ (\omterm{def:b1:topology:dense}{dichtheid}, \cref{thm:b1:reals:density}, toegepast op geneste
\omterm{prop:b1:reals:intervals}{intervallen}; of $r_n = \frac{\lfloor nx\rfloor}{n}$). \omterm{def:b1:continuity:continuous}{Continuïteit}
geeft $f(x) = \lim f(r_n) = \lim \alpha r_n = \alpha x$.
\end{solution}

\begin{solution}{exo:b1:continuity:12}
Zij $\varepsilon > 0$. Wegens de limiet in $+\infty$ is er een $A$
met $\abs{f(x) - \ell} \leq \frac\varepsilon2$ voor $x \geq A$;
bijgevolg geldt voor $x, y \geq A$ dat $\abs{f(x) - f(y)} \leq
\varepsilon$ (zonder enige eis van nabijheid).

Op het segment $\intcc{0}{A + 1}$ levert de stelling van Heine
(\cref{thm:b1:continuity:heine}) een $\delta_0 > 0$ voor deze
$\varepsilon$; stel $\delta = \min(\delta_0, 1)$.

Neem nu willekeurige $x, y \geq 0$ met $\abs{x - y} \leq \delta$,
zeg $x \leq y$. Als $y \leq A + 1$: beide liggen in het segment, en
de $\delta_0$ van Heine is van toepassing. Anders is $y > A + 1$, en
dan $x \geq y - 1 > A$: beide liggen in $\intco{A}{+\infty}$, waar
het limietargument geldt. In beide gevallen is $\abs{f(x) - f(y)}
\leq \varepsilon$: \omterm{def:b1:continuity:uniform}{uniforme continuïteit}.
\end{solution}

\begin{solution}{pb:b1:continuity:1}
\textbf{1.} Voor $n \in \N$ geldt $f(nx) = n f(x)$ met inductie
($f((n+1)x) = f(nx) + f(x)$). Verder geeft $f(0) = 2f(0)$ dat $f(0)
= 0$, en $0 = f(x - x) = f(x) + f(-x)$ dat $f$ oneven is, dus $f(nx)
= nf(x)$ voor $n \in \Z$. Voor $r = \frac pq$: $q\,f\bigl(\tfrac
pq x\bigr) = f(px) = p\,f(x)$, dus $f(rx) = r f(x)$: $f$ is
$\Q$-lineair.

\textbf{2.} De additiviteit geeft, voor alle $x$ en $h$:
\[
f(x + h) - f(x) = f(h) = f(x_0 + h) - f(x_0) .
\]
Als $h \to 0$ nadert het rechterlid tot $0$ wegens de \omterm{def:b1:continuity:continuous}{continuïteit}
in $x_0$; dus $f(x + h) \to f(x)$: \omterm{def:b1:continuity:continuous}{continuïteit} in elke $x$. Daarna
levert \cref{exo:b1:continuity:11} dat $f(x) = f(1)\,x$.

\textbf{3.} Twee \omterm{def:b1:continuity:continuous}{continue} additieve functies zijn van de vorm $cx$
en $c'x$ (vraag 2); als ze samenvallen op de \omterm{def:b1:topology:dense}{dichte} ondergroep $\Z +
\sqrt2\,\Z$ (\cref{exo:b1:reals:9}), dan is $c\,g = c'g$ voor zekere
$g \neq 0$ daarin: $c = c'$, de functies zijn gelijk. Zonder
\omterm{def:b1:continuity:continuous}{continuïteit}: $\Q$-lineariteit legt alleen de waarden in
$\Q$-combinaties vast, en $1, \sqrt2$ zijn $\Q$-onafhankelijk
($\sqrt2 \notin \Q$), dus $f(1) = 0$, $f(\sqrt2) = 1$ is consistent
en dwingt op de ondergroep
\[
f(m + n\sqrt2) = m\,f(1) + n\,f(\sqrt2) = n .
\]

\textbf{4.} Voor $t \in \intcc{0}{\ell}$ is $a + t \in
\intcc{a}{b}$, dus $f(t) = f(a + t) - f(a) \leq M - f(a) =: M'$.

\textbf{5.} Voor $t \in \intcc{0}{\ell}$ ligt ook $\ell - t$ in
$\intcc{0}{\ell}$, en de additiviteit geeft $f(t) = f(\ell) -
f(\ell - t) \geq f(\ell) - M'$. Bijgevolg is $\abs{f} \leq C$ op
$\intcc{0}{\ell}$ met $C = \max\bigl(\abs{M'}, \abs{f(\ell) -
M'}\bigr)$.

\textbf{6.} Voor $t \in \intcc{0}{\ell/n}$ is $nt \in
\intcc{0}{\ell}$ en $f(t) = \frac{f(nt)}{n}$ (vraag 1), dus
$\abs{f(t)} \leq \frac Cn$. Zij $\varepsilon > 0$ gegeven; kies $n >
\frac C\varepsilon$: voor $0 \leq h \leq \frac\ell n$ is
$\abs{f(h)} \leq \varepsilon$, en voor negatieve $h$ gebruiken we
$f(h) = -f(-h)$. Dus $f(h) \to 0 = f(0)$ als $h \to 0$:
\omterm{def:b1:continuity:continuous}{continuïteit} in $0$, dus overal (vraag 2), dus $f(x) = f(1)x$.

\textbf{7.} Als $f$ stijgend is op $\intcc{a}{b}$, dan geldt daar
$f(a) \leq f(x) \leq f(b)$: begrensd, dus lineair volgens vraag 6,
$f(x) = cx$; en $c(b - a) = f(b) - f(a) \geq 0$ dwingt $c \geq 0$.

\textbf{8.} (a)$\Rightarrow$(b)$\Rightarrow$(c): triviaal.
(c)$\Rightarrow$(a): vraag 2. (a)$\Rightarrow$(d): een lineaire
functie is overal monotoon. (d)$\Rightarrow$(e): een monotone
functie op $\intcc{a}{b}$ is daar begrensd door haar waarden in de
randpunten. (e)$\Rightarrow$(a): vragen 4--6. De vijf \omterm{def:b1:logic:statement}{uitspraken}
zijn equivalent --- de regelmaatladder.

\textbf{9.} Vraag 4 gebruikte alleen de \omterm{def:b1:reals:bounds}{bovengrens} $M$; vraag 5
\emph{leidde} de ondergrens af uit de \omterm{def:b1:reals:bounds}{bovengrens} via de spiegeling
$f(t) = f(\ell) - f(\ell - t)$; vraag 6 werkte vervolgens met
$\abs f \leq C$. Dus: additief en naar \emph{boven} begrensd op één
niet-ontaard \omterm{prop:b1:reals:intervals}{interval} impliceert al lineair. Voor een ondergrens
passen we dit toe op $-f$ (additief, naar boven begrensd). De
sterkste sport (e): een eenzijdige grens op één \omterm{prop:b1:reals:intervals}{interval} volstaat.

\textbf{10.} Als $\frac{f(x)}{x}$ voor alle $x \neq 0$ één en
dezelfde constante $c$ was, zou $f$ lineair zijn; er zijn dus $u, v
\neq 0$ met $\frac{f(u)}u \neq \frac{f(v)}v$, dat wil zeggen $u f(v)
- v f(u) \neq 0$: de determinant van de vectoren $(u, f(u))$, $(v,
f(v))$ is niet nul, en zij brengen $\R^2$ voort.

\textbf{11.} Voor $r, s \in \Q$ is $f(ru + sv) = r f(u) + s f(v)$
(tweemaal vraag 1 plus additiviteit), dus bevat de grafiek
\[
\bigl(ru + sv,\; r f(u) + s f(v)\bigr)
= r\,(u, f(u)) + s\,(v, f(v)) .
\]
Zij $(x_0, y_0)$ en $\varepsilon > 0$ gegeven: het $2 \times 2$
stelsel $a(u, f(u)) + b(v, f(v)) = (x_0, y_0)$ heeft een (unieke)
reële oplossing $(a, b)$, want de determinant is niet nul. Kies
rationale getallen $r_n \to a$, $s_n \to b$: dan geldt $r_n(u, f(u))
+ s_n(v, f(v)) \to (x_0, y_0)$ coördinaatsgewijs, en elk van die
punten ligt op de grafiek: de grafiek is \omterm{def:b1:topology:dense}{dicht} in $\R^2$.

\textbf{12.} Zij $I$ een niet-ontaard \omterm{prop:b1:reals:intervals}{interval}, $x_0$ het midden
ervan en $M$ willekeurig: de \omterm{def:b1:topology:dense}{dichtheid} levert een punt van de
grafiek op afstand minder dan $\min\bigl(\frac{\abs I}{2}, 1\bigr)$
van $(x_0, M + 1)$, dat wil zeggen een $x \in I$ met $f(x) > M$:
onbegrensd op $I$, dus (vraag 8) discontinu in elk punt en monotoon
op geen enkel \omterm{prop:b1:reals:intervals}{interval}; en voor elk doelwit $y_0$ geven punten van
de grafiek dicht bij $(x_0, y_0)$ waarden $f(x)$ willekeurig dicht
bij $y_0$ met $x \in I$: $f(I)$ is \omterm{def:b1:topology:dense}{dicht} in $\R$. Dit is vraag 8
achterstevoren gelezen: omdat alle sporten equivalent zijn, moet
een niet-lineaire additieve functie ze \emph{allemaal} missen, en
overal.

\textbf{13.} Goed gedefinieerd: uit $m + n\sqrt2 = m' + n'\sqrt2$
volgt $(n - n')\sqrt2 = m' - m \in \Z$, dus $n = n'$ (anders zou
$\sqrt2 \in \Q$) en $m = m'$. De additiviteit op $G$ is dan
coördinaat per coördinaat duidelijk. Onbegrensdheid bij $0^+$: leg
$\varepsilon \in \intoo{0}{1}$ en $N \in \N$ vast. Voor elke vaste
$n$ met $\abs n \leq N$ legt de voorwaarde $m + n\sqrt2 \in
\intoo{0}{1}$ de $m$ vast binnen een \omterm{prop:b1:reals:intervals}{interval} van lengte $1$:
hoogstens één geheel getal $m$ per $n$, dus hoogstens $2N + 1$
elementen van $G \cap \intoo{0}{1}$ hebben $\abs{f} \leq N$. Maar $G
\cap \intoo{0}{\varepsilon}$ is oneindig ($G$ is \omterm{def:b1:topology:dense}{dicht},
\cref{exo:b1:reals:9}); dus bevat die doorsnede een $g$ met
$\abs{f(g)} > N$: $f$ is onbegrensd op elke rechtse \omterm{def:b1:topology:open}{omgeving} van
$0$. $f$ uitbreiden tot een additieve functie op $\R$ betekent de
waarden coherent kiezen op een familie reële getallen die
$\Q$-lineair onafhankelijk is en $\R$ over $\Q$ voortbrengt --- een
Hamelbasis; er één produceren vergt het keuzeaxioma, en geen enkele
expliciete formule kan dat: constructief bezitten we het monster
alleen op $G$.

\textbf{14.} $f(x) = f(\frac x2 + \frac x2) = f(\frac x2)^2 \geq 0$.
Als $f(x_0) = 0$, dan is $f(x) = f(x - x_0)f(x_0) = 0$ voor alle
$x$: uitgesloten. Dus $f > 0$ en $g = \ln \circ f$ is \omterm{def:b1:continuity:continuous}{continu}
(\cref{prop:b1:continuity:algebra}) met $g(x + y) = g(x) + g(y)$:
volgens \cref{exo:b1:continuity:11} is $g(x) = cx$, dus $f(x) =
\eu^{cx}$. Omgekeerd voldoet elke $\eu^{cx}$: de \omterm{def:b1:continuity:continuous}{continue} \omterm{def:b1:structures:morphism}{morfismen}
$(\R, +) \to (\R^*, \times)$ zijn precies de exponentiëlen.

\textbf{15.} $g(u) = f(\eu^u)$ is \omterm{def:b1:continuity:continuous}{continu} en $g(u + v) = f(\eu^u
\eu^v) = g(u) + g(v)$: $g(u) = cu$, en elke $x > 0$ schrijft zich
als $x = \eu^u$ met $u = \ln x$: $f(x) = c\ln x$.

\textbf{16.} $h = \ln \circ f$ is \omterm{def:b1:continuity:continuous}{continu} op $\intoo{0}{+\infty}$
met $h(xy) = h(x) + h(y)$: volgens vraag 15 is $h(x) = c\ln x$, dus
$f(x) = \eu^{c\ln x} = x^c$.

\textbf{17.} Met $x = y = 0$: $f(0) = 2f(0) + f(0)^2$, dus $f(0)(1 +
f(0)) = 0$. Als $f(0) = -1$: door $y = 0$ te nemen, $f(x) = f(x) +
f(0) + f(x)f(0) = -1$ voor alle $x$: de constante $f \equiv -1$ (die
inderdaad aan de vergelijking voldoet). Anders is $f(0) = 0$; dan is
$h = 1 + f$ \omterm{def:b1:continuity:continuous}{continu}, $h(0) = 1$, en
\[
h(x + y) = 1 + f(x) + f(y) + f(x)f(y) = h(x)\,h(y) :
\]
volgens vraag 14 is $h(x) = \eu^{cx}$, dat wil zeggen $f(x) =
\eu^{cx} - 1$ (het geval $c = 0$ geeft $f \equiv 0$). Volledige
lijst: $f \equiv -1$ en $f(x) = \eu^{cx} - 1$, $c \in \R$.

\textbf{18.} Met $x = y = 0$: $2f(0) = 4f(0)$, dus $f(0) = 0$. Met
$x = 0$: $f(y) + f(-y) = 2f(y)$, dus $f$ is even. Met $y = x$:
$f(2x) = 4f(x)$. Inductie met $(x, y) \to (nx, x)$:
\[
f((n{+}1)x) = 2f(nx) + 2f(x) - f((n{-}1)x)
= (2n^2 + 2 - (n-1)^2) f(x) = (n+1)^2 f(x) .
\]
Vervolgens geeft $f(x) = f\bigl(q\cdot\frac xq\bigr) = q^2
f\bigl(\frac xq\bigr)$ dat $f\bigl(\frac pq x\bigr) =
\frac{p^2}{q^2}f(x)$: op $\Q$ is $f(r) = r^2 f(1)$ (de pariteit
regelt de tekens). De twee \omterm{def:b1:continuity:continuous}{continue} functies $f$ en $x \mapsto
f(1)x^2$ vallen samen op de \omterm{def:b1:topology:dense}{dichte} \omterm{def:b1:logic:sets}{verzameling} $\Q$, dus overal
(vraag 24): $f(x) = f(1)\,x^2$; elke $cx^2$ voldoet aan de
vergelijking.

\textbf{19.} $g = f - f(0)$ is \omterm{def:b1:continuity:continuous}{continu}, $g(0) = 0$, en voldoet aan
de vergelijking van Jensen (de constanten vallen weg). Met $y = 0$:
$g\bigl(\frac x2\bigr) = \frac{g(x)}{2}$. Dan geldt voor alle $x, y$
\[
\frac{g(x + y)}{2} = g\Bigl(\frac{x+y}{2}\Bigr)
= \frac{g(x) + g(y)}{2} ,
\]
dus $g$ is additief en \omterm{def:b1:continuity:continuous}{continu}: $g(x) = cx$ (vraag 2), en $f(x) = cx
+ d$ met $d = f(0)$. Alle affiene functies voldoen aan Jensen: de
lijst is volledig.

\textbf{20.} Inductie op $n$. Voor $n = 0$: $\lambda \in \{0, 1\}$,
triviaal. Neem aan dat de ongelijkheid geldt voor alle gewichten
$\frac{k}{2^n}$. Een gewicht $\lambda = \frac{k}{2^{n+1}}$ met $k$
even herleidt zich tot niveau $n$; voor $k = 2j + 1$ is $\lambda$
het midden van $\lambda_1 = \frac{j}{2^n}$ en $\lambda_2 =
\frac{j+1}{2^n}$. Met $z_i = \lambda_i x + (1 - \lambda_i)y$ is
$\lambda x + (1-\lambda)y = \frac{z_1 + z_2}{2}$, dus
\[
f\bigl(\lambda x + (1{-}\lambda)y\bigr)
\leq \frac{f(z_1) + f(z_2)}{2}
\leq \frac{(\lambda_1 + \lambda_2)f(x) + (2 - \lambda_1 -
\lambda_2)f(y)}{2}
= \lambda f(x) + (1 - \lambda)f(y) .
\]

\textbf{21.} Leg $x, y$ vast. De \omterm{def:b1:logic:map}{afbeeldingen} $\lambda \mapsto
f(\lambda x + (1 - \lambda)y)$ en $\lambda \mapsto \lambda f(x) + (1
- \lambda)f(y)$ zijn \omterm{def:b1:continuity:continuous}{continu} op $\intcc{0}{1}$ (samenstelling en
algebra, \cref{prop:b1:continuity:algebra}). De ongelijkheid geldt
op de dyadische gewichten, die \omterm{def:b1:topology:dense}{dicht} liggen in $\intcc{0}{1}$
(\cref{exo:b1:reals:8}); voor een willekeurige $\lambda$ nemen we
dyadische $\lambda_n \to \lambda$ en gaan we over tot de limiet
(\cref{thm:b1:continuity:seqchar} en \cref{thm:b1:seq:order}): de
convexiteitsongelijkheid geldt voor elke $\lambda \in
\intcc{0}{1}$ --- middenpuntsconvexiteit plus \omterm{def:b1:continuity:continuous}{continuïteit} is
convexiteit (het begrip uit \cref{ch:b1:derivative}).

\textbf{22.} Een niet-lineaire additieve $f$ voldoet exact aan
$f\bigl(\frac{x+y}{2}\bigr) = \frac{f(x) + f(y)}{2}$ (vraag 1 met $r
= \frac12$, daarna additiviteit): zij is middenpuntsconvex, zelfs
middenpuntsaffien. Als zij op een \omterm{prop:b1:reals:intervals}{interval} $\intcc{x}{y}$ aan de
volledige convexiteitsongelijkheid voldeed, dan gold voor $\lambda
\in \intcc{0}{1}$: $f(\lambda x + (1-\lambda)y) \leq \max(f(x),
f(y))$ --- naar boven begrensd op een niet-ontaard \omterm{prop:b1:reals:intervals}{interval}, dus
lineair volgens vraag 9: tegenspraak. De \omterm{def:b1:continuity:continuous}{continuïteit} in vraag 21 is
dus geen luxe: zonder haar beheerst de middenpuntsconvexiteit alleen
het aftelbare dyadische skelet, en het continuüm daartussen loopt
uit de hand.

\textbf{23.} $p(x) = \frac{x^2}{2}$ voldoet aan $p(x+y) = p(x) +
p(y) + xy$. Is $f$ een willekeurige \omterm{def:b1:continuity:continuous}{continue} oplossing, dan is $g =
f - p$ \omterm{def:b1:continuity:continuous}{continu} en additief, dus $g(x) = cx$:
\[
f(x) = \frac{x^2}{2} + cx , \qquad c \in \R ,
\]
en elk van deze functies is een oplossing: de lijst is volledig.

\textbf{24.} Algemeen principe: zijn $u, v$ \omterm{def:b1:continuity:continuous}{continu} en vallen ze
samen op een \omterm{def:b1:topology:dense}{dichte} $D \subseteq \R$, kies dan voor $x \in \R$
punten $d_n \in D$ met $d_n \to x$
(\cref{prop:b1:topology:closureprops}); dan is $u(x) = \lim u(d_n) =
\lim v(d_n) = v(x)$. Zij nu $f$ \omterm{def:b1:continuity:continuous}{continu} en additief op de \omterm{def:b1:topology:dense}{dichte}
ondergroep $G$. Leg $x, y \in \R$ vast en neem $g_n \to x$, $g'_n
\to y$ met $g_n, g'_n \in G$; dan geldt $g_n + g'_n \to x + y$ en,
wegens de rijcontinuïteit in $x + y$, $x$ en $y$,
\[
f(x + y) = \lim f(g_n + g'_n) = \lim\bigl(f(g_n) +
f(g'_n)\bigr) = f(x) + f(y) :
\]
$f$ is additief op heel $\R$ (en dus lineair, volgens vraag 2).

\textbf{25.} (i) Voor additieve $f$: lineair $\iff$ \omterm{def:b1:continuity:continuous}{continu} $\iff$
\omterm{def:b1:continuity:continuous}{continu} in één punt $\iff$ monotoon op een of ander \omterm{prop:b1:reals:intervals}{interval} $\iff$
begrensd (zelfs eenzijdig) op een of ander \omterm{prop:b1:reals:intervals}{interval}. (ii) De
\omterm{def:b1:topology:dense}{dichtheid} van de grafiek in het vlak laat zien dat het tekort geen
plaatselijk gebrek is maar een globale explosie: boven elk
deelinterval smeren de waarden zich uit over heel $\R$, zodat elke
regelmaatseigenschap overal tegelijk faalt. (iii) De vangst: $cx$,
$\eu^{cx}$, $c\ln x$, $x^c$, $cx^2$ en $cx + d$ --- zes
karakteriseringen, één methode: breng de vergelijking terug tot die
van Cauchy, bewijs het $\Q$-skelet met inductie, til het op tot $\R$
met \omterm{def:b1:topology:dense}{dichtheid} plus \omterm{def:b1:continuity:continuous}{continuïteit}. (iv) De \omterm{def:b1:topology:dense}{dichtheid} van $\Q$ (of van
de dyadische getallen) droeg de opwaarderingen in
\cref{exo:b1:continuity:11} en in vraag 21; ze droeg niets in vraag
13, want zonder \omterm{def:b1:continuity:continuous}{continuïteit} planten waarden zich niet voort van een
\omterm{def:b1:topology:dense}{dichte} \omterm{def:b1:logic:sets}{verzameling} naar haar \omterm{def:b1:topology:closure}{afsluiting} --- \omterm{def:b1:topology:dense}{dichtheid} draagt
informatie alleen langs \omterm{def:b1:continuity:continuous}{continuïteit}.
\end{solution}
